% Original analytic teaching diagram, 2026-11-09.
\definecolor{VB}{HTML}{7188A3}\definecolor{VBF}{HTML}{E5EAF0}
\definecolor{VG}{HTML}{829C8B}\definecolor{VGF}{HTML}{E8EEE9}
\definecolor{VY}{HTML}{B4A16F}\definecolor{VYF}{HTML}{F1ECDD}
\definecolor{VR}{HTML}{A7787D}\definecolor{VRF}{HTML}{F0E5E6}
\definecolor{VP}{HTML}{9588A5}\definecolor{VPF}{HTML}{ECE8F0}
\definecolor{VI}{HTML}{3B4252}\definecolor{VM}{HTML}{6B7280}\definecolor{VL}{HTML}{4C4D4F}
\begin{tikzpicture}[x=1mm,y=1mm,>=Stealth,font=\figurefont\fontsize{8}{10}\selectfont,
lab/.style={text=FigureMuted,font=\figurefont\fontsize{7.5}{9}\selectfont},
box/.style={draw=VB,fill=VBF,rounded corners=1pt,text=VI,align=center,minimum height=10mm,inner sep=1mm,text width=23mm},
dual/.style={draw=VP,fill=VPF,rounded corners=1pt,text=VI,align=center,minimum height=10mm,inner sep=1mm,text width=23mm},
ints/.style={draw=VY,fill=VYF,rounded corners=1pt,text=VI,align=center,minimum height=10mm,inner sep=1mm,text width=23mm},
arr/.style={->,draw=VL,line width=.7pt},note/.style={text=VI,font=\figurefont\fontsize{7.5}{9}\selectfont,fill=white,inner sep=.8pt}]
\node[box] (a) at (0,0) {线性模型\\（等价变换）};
\node[box] (b) at (36,0) {标准型\\满行秩};
\node[box] (c) at (72,0) {基与基本可行解\\（极点）};
\node[box] (d) at (72,-22) {最优基本可行解};
\node[dual] (e) at (36,-22) {对偶问题\\弱对偶};
\node[dual] (f) at (36,-44) {强对偶\\零间隙证书};
\node[ints] (g) at (0,-22) {整数问题};
\node[ints] (h) at (0,-44) {线性松弛\\下界};
\node[ints] (i) at (72,-44) {全幺模\\整数顶点};
\draw[arr] (a)--(b); \draw[arr] (b)--(c); \draw[arr] (c)--node[note,right]{可行且有下界}(d);
\draw[arr] (b)--(e); \draw[arr] (e)--node[note,right]{可行且有下界}(f);
\draw[arr] (g)--(h); \draw[arr] (h.east)--(f.west);
\draw[arr] (d)--(i);
\end{tikzpicture}
